% GENERATED FILE. Edit canonical lessons, then run npm run generate:books.
% canonical-source-hash: fnv1a64:e4e070aeb7ff8b4d
% canonical-lessons: SA-C170-karyam, SA-C170-karma, SA-C170-udyogah, SA-C170-vrttih, SA-C170-seva

\chapter{Work, Action, Occupation, Livelihood, Service}
\label{ch:sa-work-action-occupation-livelihood-service}
\hlchaptermodality{170}

\begin{chapteropening}
\textbf{By the end of this chapter:} \emph{I can say work, an action, an occupation, a livelihood and service.}

Complete the last lesson of chapter 170: I can say work, an action, an occupation, a livelihood and service.
\end{chapteropening}

\section[\texorpdfstring{\sk{कार्यम्} (kāryam)}{कार्यम् (kāryam)}]{\sk{कार्यम्} (kāryam) — work, a task}

\label{lesson:SA-C170-karyam}



\begin{quote}
Before the new one: say the Sanskrit for a dictionary, then the Sanskrit for writing.
\end{quote}

\subsection*{You'll want to know: \sk{कार्यम्}}
\textbf{\sk{कार्यम्}} — \emph{kāryam} — \textquotedblleft{}work, a task\textquotedblright{}.

\begin{grammarlens}[title={What you've built}]
One more word for this chapter.
\end{grammarlens}

\subsection*{Your turn}
\begin{itemize}
\raggedright
  \item \emph{Say it:} \emph{kāryam}
  \item \emph{Say it:} \emph{kāryam}, once more
  \item \emph{Say it:} say \emph{kāryam}
  \item \emph{Recall:} say the Sanskrit for a dictionary, then the Sanskrit for writing, then say \emph{kāryam} again
\end{itemize}

\begin{tcolorbox}[breakable,colback=teal!4,colframe=teal!35!black,title={Before you move on}]
What does \sk{कार्यम्} mean? (\textquotedblleft{}Work, a task\textquotedblright{}.) Say it once more.
\end{tcolorbox}
\section[\texorpdfstring{\sk{कर्म} (karma)}{कर्म (karma)}]{\sk{कर्म} (karma) — an action, work}

\label{lesson:SA-C170-karma}



\begin{quote}
Before the new one: say the Sanskrit for writing, then the Sanskrit for work.
\end{quote}

\subsection*{You'll want to know: \sk{कर्म}}
\textbf{\sk{कर्म}} — \emph{karma} — \textquotedblleft{}an action, work\textquotedblright{}.

\begin{grammarlens}[title={What you've built}]
One more word for this chapter.
\end{grammarlens}

\subsection*{Your turn}
\begin{itemize}
\raggedright
  \item \emph{Say it:} \emph{karma}
  \item \emph{Say it:} \emph{karma}, once more
  \item \emph{Say it:} say \emph{karma}
  \item \emph{Recall:} say the Sanskrit for writing, then the Sanskrit for work, then say \emph{karma} again
\end{itemize}

\begin{tcolorbox}[breakable,colback=teal!4,colframe=teal!35!black,title={Before you move on}]
What does \sk{कर्म} mean? (\textquotedblleft{}An action, work\textquotedblright{}.) Say it once more.
\end{tcolorbox}
\section[\texorpdfstring{\sk{उद्योगः} (udyogaḥ)}{उद्योगः (udyogaḥ)}]{\sk{उद्योगः} (udyogaḥ) — an occupation, industry}

\label{lesson:SA-C170-udyogah}



\begin{quote}
Before the new one: say the Sanskrit for work, then the Sanskrit for an action.
\end{quote}

\subsection*{You'll want to know: \sk{उद्योगः}}
\textbf{\sk{उद्योगः}} — \emph{udyogaḥ} — \textquotedblleft{}an occupation, industry\textquotedblright{}.

\begin{grammarlens}[title={What you've built}]
One more word for this chapter.
\end{grammarlens}

\subsection*{Your turn}
\begin{itemize}
\raggedright
  \item \emph{Say it:} \emph{udyogaḥ}
  \item \emph{Say it:} \emph{udyogaḥ}, once more
  \item \emph{Say it:} say \emph{udyogaḥ}
  \item \emph{Recall:} say the Sanskrit for work, then the Sanskrit for an action, then say \emph{udyogaḥ} again
\end{itemize}

\begin{tcolorbox}[breakable,colback=teal!4,colframe=teal!35!black,title={Before you move on}]
What does \sk{उद्योगः} mean? (\textquotedblleft{}An occupation, industry\textquotedblright{}.) Say it once more.
\end{tcolorbox}
\section[\texorpdfstring{\sk{वृत्तिः} (vṛttiḥ)}{वृत्तिः (vṛttiḥ)}]{\sk{वृत्तिः} (vṛttiḥ) — a livelihood}

\label{lesson:SA-C170-vrttih}



\begin{quote}
Before the new one: say the Sanskrit for an action, then the Sanskrit for an occupation.
\end{quote}

\subsection*{You'll want to know: \sk{वृत्तिः}}
\textbf{\sk{वृत्तिः}} — \emph{vṛttiḥ} — \textquotedblleft{}a livelihood\textquotedblright{}.

\begin{grammarlens}[title={What you've built}]
One more word for this chapter.
\end{grammarlens}

\subsection*{Your turn}
\begin{itemize}
\raggedright
  \item \emph{Say it:} \emph{vṛttiḥ}
  \item \emph{Say it:} \emph{vṛttiḥ}, once more
  \item \emph{Say it:} say \emph{vṛttiḥ}
  \item \emph{Recall:} say the Sanskrit for an action, then the Sanskrit for an occupation, then say \emph{vṛttiḥ} again
\end{itemize}

\begin{tcolorbox}[breakable,colback=teal!4,colframe=teal!35!black,title={Before you move on}]
What does \sk{वृत्तिः} mean? (\textquotedblleft{}A livelihood\textquotedblright{}.) Say it once more.
\end{tcolorbox}
\section[\texorpdfstring{\sk{सेवा} (sevā)}{सेवा (sevā)}]{\sk{सेवा} (sevā) — service}

\label{lesson:SA-C170-seva}



\begin{quote}
Before the new one: say the Sanskrit for an occupation, then the Sanskrit for a livelihood.
\end{quote}

\subsection*{You'll want to know: \sk{सेवा}}
\textbf{\sk{सेवा}} — \emph{sevā} — \textquotedblleft{}service\textquotedblright{}.

\begin{grammarlens}[title={What you've built}]
One more word for this chapter.
\end{grammarlens}

\subsection*{Your turn}
\begin{itemize}
\raggedright
  \item \emph{Say it:} \emph{sevā}
  \item \emph{Say it:} \emph{sevā}, once more
  \item \emph{Say it:} say \emph{sevā}
  \item \emph{Recall:} say the Sanskrit for an occupation, then the Sanskrit for a livelihood, then say \emph{sevā} again
\end{itemize}

\begin{tcolorbox}[breakable,colback=teal!4,colframe=teal!35!black,title={Before you move on}]
What does \sk{सेवा} mean? (\textquotedblleft{}Service\textquotedblright{}.) Say it once more.
\end{tcolorbox}
